\documentclass[11pt]{article}

% Use a NeurIPS-like format
\usepackage[utf8]{inputenc}
\usepackage[T1]{fontenc}
\usepackage{times}
\usepackage[margin=1in]{geometry}
\usepackage{amsmath,amssymb,amsfonts}
\usepackage{booktabs}
\usepackage{graphicx}
\usepackage{hyperref}
\usepackage{algorithm}
\usepackage{algorithmic}
\usepackage{natbib}
\usepackage{xcolor}
\usepackage{microtype}
\usepackage{subcaption}
\usepackage{multirow}

\newcommand{\E}{\mathbb{E}}
\newcommand{\Var}{\mathrm{Var}}
\newcommand{\R}{\mathbb{R}}
\newcommand{\N}{\mathcal{N}}
\newcommand{\CS}{\mathrm{CS}}
\newcommand{\CI}{\mathrm{CI}}
\newcommand{\iid}{\mathrm{i.i.d.}}

\title{Confidence Sequences for Markov Chain Monte Carlo:\\ Anytime-Valid Estimation with Sequential Guarantees}

\author{
  Anonymous Author(s)
}

\date{}

\begin{document}
\maketitle

% ============================================================================
% ABSTRACT
% ============================================================================
\begin{abstract}
Markov chain Monte Carlo (MCMC) practitioners routinely monitor estimates and stop sampling based on heuristic criteria, yet existing confidence intervals are only valid at pre-specified sample sizes. This creates a fundamental gap: the sequential monitoring commonly used in practice violates the fixed-sample assumptions underlying standard CLT-based intervals. We develop the first confidence sequences for MCMC estimators---time-uniform confidence intervals that remain valid at all stopping times simultaneously. We propose two complementary methods: Batch Means Confidence Sequences (BMCS), which partition chains into adaptive batches and apply empirical Bernstein bounds, and Coupling-Based Confidence Sequences (CBCS), which leverage unbiased MCMC to generate exactly independent estimators. In experiments on Gaussian targets up to $d=50$ dimensions and Bayesian logistic regression on three UCI datasets, BMCS maintains $\geq 93.7\%$ time-uniform coverage (nominal 95\%) while CLT intervals degrade to as low as 28.3\% under sequential monitoring. BMCS adds less than 20\% computational overhead to the base MCMC chain. Our results demonstrate that anytime-valid inference for MCMC is practical and provides substantially stronger guarantees than current practice, though at the cost of wider intervals---approximately $1.9\times$ the CLT width.
\end{abstract}

% ============================================================================
% INTRODUCTION
% ============================================================================
\section{Introduction}

Markov chain Monte Carlo (MCMC) is the workhorse of Bayesian computation, enabling approximate inference in complex probabilistic models by constructing Markov chains whose stationary distributions match the target posterior. A central practical challenge in MCMC is deciding \emph{when to stop}: practitioners must determine when the chain has run long enough that the resulting estimates are sufficiently accurate.

The standard approach constructs confidence intervals using the Markov chain central limit theorem (CLT), estimating the asymptotic variance via batch means or spectral methods \citep{flegal2010batch,vats2019multivariate}. Fixed-width stopping rules (FWSR) terminate sampling when the confidence interval half-width falls below a threshold \citep{vats2019multivariate,flegal2024implementing}. However, these intervals are only valid at a \emph{pre-specified} sample size $n$. When practitioners monitor confidence intervals sequentially and stop as soon as they are narrow enough---as is universal in practice---the resulting coverage guarantees are violated due to the well-known sequential testing problem.

In the sequential analysis literature, \emph{confidence sequences} (CS) provide time-uniform confidence intervals that are valid at \emph{all} stopping times simultaneously \citep{howard2021time,ramdas2023game}. A $(1-\alpha)$ confidence sequence satisfies
\begin{equation}
\Pr\!\left(\mu \in \CS_t \;\text{for all}\; t \geq 1\right) \geq 1 - \alpha,
\label{eq:cs_def}
\end{equation}
which is strictly stronger than a fixed-time guarantee $\Pr(\mu \in \CI_n) \geq 1-\alpha$ for a single pre-specified~$n$. Despite the maturity of both literatures, \emph{no existing work develops confidence sequences specifically for MCMC estimators}. This gap is surprising given that MCMC is inherently sequential.

The key challenge is that MCMC samples are dependent, precluding direct application of i.i.d.\ confidence sequence constructions. Our insight is that two established techniques can bridge this gap: (1) \textbf{batch means} convert dependent MCMC samples into approximately independent averages, and (2) \textbf{coupling-based debiasing} \citep{jacob2020unbiased} produces exactly independent and unbiased estimators from coupled chains.

\paragraph{Contributions.} Our contributions are:
\begin{itemize}
  \item We develop the first confidence sequences for MCMC estimators, providing anytime-valid coverage guarantees that hold at all stopping times simultaneously.
  \item We propose two methods---BMCS (batch means) and CBCS (coupling-based)---with complementary tradeoffs between validity, width, and computational cost.
  \item We provide comprehensive experiments on Gaussian targets ($d \in \{1, 10, 50\}$), Bayesian logistic regression on three UCI datasets, and multimodal targets, demonstrating that BMCS maintains $\geq 93.7\%$ time-uniform coverage while CLT intervals degrade to as low as 28.3\% under sequential monitoring.
  \item We characterize the practical cost of anytime validity: BMCS intervals are approximately $1.9\times$ wider than CLT intervals but add less than 20\% computational overhead.
\end{itemize}

The remainder of this paper is organized as follows. Section~\ref{sec:related} reviews related work. Section~\ref{sec:method} describes our proposed methods. Section~\ref{sec:experiments} presents experimental results, and Section~\ref{sec:discussion} discusses limitations and conclusions.

% ============================================================================
% RELATED WORK
% ============================================================================
\section{Related Work}
\label{sec:related}

\paragraph{MCMC output analysis.}
The foundational work on MCMC confidence intervals uses the Markov chain CLT and asymptotic variance estimation via batch means \citep{flegal2010batch}, spectral methods \citep{geyer1992practical}, or the initial sequence estimator \citep{geyer2011introduction}. \citet{vats2019multivariate} extend this to multivariate settings and propose effective sample size and fixed-width stopping rules. \citet{vats2022lugsail} introduce lugsail lag windows to reduce negative bias in variance estimators. \citet{flegal2024implementing} provide a comprehensive practical workflow. All existing MCMC confidence intervals are valid only at a pre-specified sample size. Our confidence sequences are valid at \emph{all} stopping times, enabling principled sequential monitoring.

\paragraph{Confidence sequences and anytime-valid inference.}
\citet{howard2021time} develop the modern theory of confidence sequences for i.i.d.\ data, achieving nonasymptotic, nonparametric, time-uniform coverage. \citet{ramdas2023game} survey the game-theoretic foundations, connecting confidence sequences to e-values and test martingales. \citet{wang2023catoni} extend to heavy-tailed distributions using Catoni's M-estimator. Existing confidence sequences assume i.i.d.\ or martingale data. Our work extends these constructions to the dependent, non-stationary setting of MCMC, requiring new technical tools including batch means approximation and coupling-based debiasing.

\paragraph{Coupling-based MCMC methods.}
\citet{jacob2020unbiased} introduce unbiased MCMC with couplings, using meeting times of coupled chains to construct exactly unbiased estimators. \citet{corenflos2025coupling} use couplings for f-divergence-based convergence diagnostics. \citet{martinez2025sequential} develop sequential goodness-of-fit tests using kernel Stein discrepancy. Jacob et al.\ provide unbiased point estimators but not confidence sequences; Martinez-Taboada \& Ramdas address testing but not estimation. Our work provides the first \emph{estimation} confidence sequences for MCMC.

\paragraph{MCMC stopping rules.}
\citet{robertson2021assessing} and \citet{vats2019multivariate} propose fixed-width stopping rules that terminate when the CLT-based confidence interval is narrow enough. These rules have desirable asymptotic properties but lack finite-sample sequential validity. Our confidence sequences maintain valid coverage at \emph{every} intermediate time point, providing strictly stronger guarantees.

% ============================================================================
% METHOD
% ============================================================================
\section{Method}
\label{sec:method}

\subsection{Problem Setup}

Let $\{X_1, X_2, \ldots\}$ be an MCMC chain targeting distribution~$\pi$ on $\R^d$, and let $f:\R^d \to \R$ be a function whose expectation $\mu = \E_\pi[f(X)]$ we wish to estimate. The running average is $\bar{f}_n = \frac{1}{n}\sum_{i=1}^n f(X_i)$. Under regularity conditions (geometric ergodicity), the Markov chain CLT gives
\begin{equation}
\sqrt{n}(\bar{f}_n - \mu) \xrightarrow{d} \N(0, \sigma^2_\infty),
\label{eq:mcmc_clt}
\end{equation}
where $\sigma^2_\infty = \sum_{k=-\infty}^{\infty} \mathrm{Cov}(f(X_0), f(X_k))$ is the asymptotic variance. The standard CLT-based confidence interval at sample size~$n$ is
\begin{equation}
\CI_n = \left[\bar{f}_n \pm z_{\alpha/2} \sqrt{\hat{\sigma}^2_\infty / n}\right],
\label{eq:clt_ci}
\end{equation}
where $\hat{\sigma}^2_\infty$ is estimated via batch means. This interval is valid only at a single, pre-specified~$n$. Our goal is to construct a confidence \emph{sequence} $(\CS_t)_{t \geq 1}$ satisfying~\eqref{eq:cs_def}.

\subsection{Method 1: Batch Means Confidence Sequences (BMCS)}

\paragraph{Core idea.} Partition the MCMC chain into non-overlapping batches and apply confidence sequence constructions to the resulting batch means, which are approximately independent for sufficiently large batch sizes.

Given the chain $\{X_1, \ldots, X_n\}$, define batch means of size~$b$:
\begin{equation}
Y_j = \frac{1}{b}\sum_{i=(j-1)b+1}^{jb} f(X_i), \quad j = 1, 2, \ldots, k = \lfloor n/b \rfloor.
\label{eq:batch_means}
\end{equation}
Under geometric ergodicity, for large~$b$, the batch means $Y_1, \ldots, Y_k$ are approximately independent with $\E[Y_j] \approx \mu$ and $\Var(Y_j) \approx \sigma^2_\infty / b$.

\paragraph{Adaptive batch size.} We use an adaptive batch size $b_k = \max(50, \lceil c \cdot \log(k+1) \rceil)$ for a constant~$c$ calibrated from an initial mixing time estimate. This ensures enough batches for tight confidence sequences while controlling residual dependence.

\paragraph{Confidence sequence construction.} After $k$ batches, with running mean $\bar{Y}_k = \frac{1}{k}\sum_{j=1}^k Y_j$ and empirical variance $\hat{V}_k = \frac{1}{k-1}\sum_{j=1}^k (Y_j - \bar{Y}_k)^2$, we construct the empirical Bernstein confidence sequence \citep{howard2021time}:
\begin{equation}
\CS_k^{\text{BMCS}} = \left[\bar{Y}_k \pm \sqrt{\frac{2\hat{V}_k \log(\log(2k)/\alpha)}{k}} + \frac{3B\log(\log(2k)/\alpha)}{k} + \delta_k\right],
\label{eq:bmcs}
\end{equation}
where $B$ is a bound on $|Y_j - \mu|$, and $\delta_k = C\rho^{b_k}$ is a correction term for residual batch dependence with $\rho < 1$ being the geometric mixing rate. The first term adapts to the empirical variance, the second provides a worst-case bound for small~$k$, and the third controls approximation error from dependence.

\paragraph{Properties.} The BMCS width converges at rate $O(\sqrt{\log\log k / k})$ in the number of batches, translating to $O(\sqrt{\log\log(n/b) / (n/b)})$ in total samples---near-optimal up to the $\log\log$ factor.

\subsection{Method 2: Coupling-Based Confidence Sequences (CBCS)}

\paragraph{Core idea.} Use the unbiased MCMC framework of \citet{jacob2020unbiased} to generate exactly independent and unbiased estimators, then apply standard i.i.d.\ confidence sequences directly.

The method runs coupled pairs of chains $(X_t, Y_t)$ with a maximal coupling kernel. When the chains meet at time~$\tau$ (i.e., $X_\tau = Y_\tau$), unbiased estimators are constructed as:
\begin{equation}
H_k = f(X_k) + \sum_{t=k+1}^{\tau-1} [f(X_t) - f(Y_{t-1})],
\label{eq:unbiased_estimator}
\end{equation}
with burn-in parameter $k = \lfloor \tau/2 \rfloor$. These estimators satisfy $\E[H_k] = \mu$ exactly and can be generated independently by running multiple coupled chain pairs.

Given $m$ independent unbiased estimators $H_1, \ldots, H_m$, we apply the sub-Gaussian confidence sequence:
\begin{equation}
\CS_m^{\text{CBCS}} = \left[\bar{H}_m \pm \sqrt{\frac{2\hat{\sigma}^2_H \log(\log(2m)/\alpha)}{m}}\right],
\label{eq:cbcs}
\end{equation}
where $\hat{\sigma}^2_H$ is the empirical variance of the unbiased estimators.

\paragraph{Properties.} CBCS provides \emph{exact} independence, eliminating the approximation error inherent in BMCS. However, it incurs computational overhead from coupling (each estimator requires running two coupled chains until meeting) and higher variance per estimator due to the telescopic correction in~\eqref{eq:unbiased_estimator}.

\subsection{Method 3: Martingale Decomposition Confidence Sequences (MDCS)}

Under suitable regularity conditions, the MCMC partial sum admits a Poisson equation decomposition:
\begin{equation}
S_n - n\mu = M_n + R_n,
\label{eq:poisson}
\end{equation}
where $M_n$ is a martingale and $R_n = g(X_n) - g(X_0)$ is a bounded remainder with $g$ being the Poisson equation solution. MDCS applies martingale confidence sequences to $M_n$ and bounds $R_n$ separately. This approach is the most direct but requires estimating the Poisson equation solution norm, which we approximate online via kernel-based regression. As our experiments show, this estimation is challenging in practice, leading to less reliable coverage in high dimensions.

\subsection{Practical Comparison}

Table~\ref{tab:method_comparison} summarizes the tradeoffs.

\begin{table}[t]
\caption{Comparison of proposed confidence sequence methods for MCMC.}
\label{tab:method_comparison}
\centering
\begin{tabular}{lcccc}
\toprule
Method & Independence & Width & Overhead & Assumptions \\
\midrule
BMCS & Approximate & Moderate & Low ($<$20\%) & Geometric ergodicity \\
CBCS & Exact & Variable & High (coupling) & Faithful coupling \\
MDCS & Via decomposition & Variable & Moderate & Drift + minorization \\
\bottomrule
\end{tabular}
\end{table}

Based on our experimental findings, we recommend \textbf{BMCS as the default practical choice}: it provides reliable coverage with minimal overhead. CBCS is suitable when exact validity is paramount and coupling is efficient (low dimensions with fast meeting times). MDCS, while theoretically appealing, requires further development for reliable practical use.

% ============================================================================
% EXPERIMENTS
% ============================================================================
\section{Experiments}
\label{sec:experiments}

All experiments run on CPU. We implement all methods in Python using NumPy/SciPy, with MCMC samplers written from scratch for full control. We use $\alpha = 0.05$ (nominal 95\% coverage) throughout and report results averaged over 3 random seeds (42, 123, 456).

\subsection{Experiment 1: Coverage Verification on Gaussian Targets}

\paragraph{Setup.} We target $\N(0, I_d)$ for $d \in \{1, 10, 50\}$ with known mean $\mu = 0$, enabling exact coverage assessment. We run 200 independent Random Walk Metropolis-Hastings (RWMH) replicates per seed (600 total) with chains of length $n = 50{,}000$. At monitoring times $t \in \{100, 200, 500, 1000, 2000, 5000, 10000, 20000, 50000\}$, we check whether $\mu = 0$ falls within each method's interval and compute time-uniform coverage as $\min_t \text{coverage}(t)$.

\paragraph{Results.} Table~\ref{tab:coverage} presents the main coverage results. The key findings are:

\begin{table}[t]
\caption{Time-uniform coverage (\%) and final-time coverage (\%) on Gaussian targets. Nominal level is 95\%. $\uparrow$ means higher is better. Best time-uniform coverage in \textbf{bold}. 600 replicates per cell.}
\label{tab:coverage}
\centering
\begin{tabular}{lcccccc}
\toprule
& \multicolumn{2}{c}{$d=1$} & \multicolumn{2}{c}{$d=10$} & \multicolumn{2}{c}{$d=50$} \\
\cmidrule(lr){2-3} \cmidrule(lr){4-5} \cmidrule(lr){6-7}
Method & Unif.\ $\uparrow$ & Final $\uparrow$ & Unif.\ $\uparrow$ & Final $\uparrow$ & Unif.\ $\uparrow$ & Final $\uparrow$ \\
\midrule
CLT CI     & 88.2 & 93.3 & 62.5 & 95.0 & 28.3 & 88.5 \\
BMCS       & \textbf{98.2} & \textbf{100.0} & \textbf{98.7} & \textbf{100.0} & \textbf{93.7} & \textbf{100.0} \\
CBCS       & 99.7 & 99.7 & 99.8 & 99.8 & \textbf{100.0} & \textbf{100.0} \\
MDCS       & 89.8 & 100.0 & 89.0 & 100.0 & 65.7 & 100.0 \\
\bottomrule
\end{tabular}
\end{table}

\textbf{CLT intervals fail under sequential monitoring.} While CLT intervals achieve near-nominal coverage at the final time $n=50{,}000$, their time-uniform coverage degrades dramatically: from 88.2\% at $d=1$ to just 28.3\% at $d=50$. This confirms that the standard practice of monitoring CLT intervals and stopping early leads to severe coverage violations.

\textbf{BMCS maintains valid coverage.} BMCS achieves $\geq 93.7\%$ time-uniform coverage across all dimensions, close to the nominal 95\%. The slight shortfall at $d=50$ reflects the challenge of batch means approximation in higher dimensions where mixing is slower.

\textbf{CBCS achieves excellent coverage} ($\geq 99.7\%$) due to exact independence of the unbiased estimators. However, as we discuss below, this comes at significant computational cost.

\textbf{MDCS has unreliable coverage.} The martingale decomposition approach achieves only 89.0\% at $d=10$ and 65.7\% at $d=50$, falling well below nominal. The difficulty of accurately estimating the Poisson equation solution norm in practice limits this method's reliability.

Figure~\ref{fig:coverage} shows coverage trajectories over time. BMCS and CBCS consistently stay above the 95\% line (with BMCS slightly dipping early in $d=50$), while CLT coverage starts low and slowly climbs.

\begin{figure}[t]
\centering
\includegraphics[width=\textwidth]{figures/fig1_coverage_trajectories.pdf}
\caption{Coverage trajectories over sample size for $d \in \{1, 10, 50\}$. Dashed horizontal line marks 95\% nominal level. CLT intervals (gray) exhibit severe undercoverage under sequential monitoring, especially in higher dimensions. BMCS (blue) and CBCS (orange) maintain coverage above the nominal level. MDCS (green) is unreliable, particularly for large $d$.}
\label{fig:coverage}
\end{figure}

\subsection{Width Convergence}

Figure~\ref{fig:width} shows the half-width convergence of each method for $d=10$. All methods converge at approximately the $O(1/\sqrt{n})$ rate, with BMCS width approximately $1.86\times$ the CLT width at $n=50{,}000$ ($0.089$ vs.\ $0.048$). This ``price of anytime validity'' is the cost of maintaining coverage at all stopping times rather than just one.

\begin{figure}[t]
\centering
\includegraphics[width=0.55\textwidth]{figures/fig2_width_trajectories.pdf}
\caption{Confidence sequence half-width vs.\ sample size for $d=10$ on log-log axes. The dashed line shows the optimal $O(1/\sqrt{n})$ rate. BMCS (blue) tracks close to the CLT CI (gray) with an approximately $1.9\times$ multiplicative overhead. MDCS (green) produces wider intervals due to conservative Poisson equation norm estimates.}
\label{fig:width}
\end{figure}

\subsection{Experiment 2: Bayesian Logistic Regression}

\paragraph{Setup.} We evaluate on Bayesian logistic regression with $\beta \sim \N(0, 10I)$ prior on three UCI datasets: German Credit ($n=1000$, $p=21$), Ionosphere ($n=351$, $p=34$), and Pima Indians Diabetes ($n=768$, $p=9$). We run RWMH chains of length $100{,}000$ with 3 seeds, estimating posterior means of regression coefficients. Reference values are computed from very long chains. Coverage is assessed across all coordinates.

\paragraph{Results.} Table~\ref{tab:logistic} presents the results. BMCS achieves 100\% coverage on German Credit and Pima, with 80\% on Ionosphere (where the higher dimensionality of $p=34$ makes mixing more challenging). CLT intervals show poor sequential coverage: 46.7\% on German Credit and 40.0\% on Ionosphere. CBCS largely fails in this setting---coupling does not converge within the computational budget for German Credit and Ionosphere, and produces very wide intervals ($0.77$) with 0\% coordinate coverage on Pima. This confirms that CBCS is impractical for moderate-dimensional problems.

\begin{table}[t]
\caption{Bayesian logistic regression results. Coverage (\%) is time-uniform across coordinates. Width is the mean half-width at $n=100{,}000$. $\uparrow$ means higher is better for coverage, $\downarrow$ for width. ``---'' indicates coupling did not converge.}
\label{tab:logistic}
\centering
\begin{tabular}{llccc}
\toprule
Dataset & Method & Coverage (\%) $\uparrow$ & Width $\downarrow$ \\
\midrule
\multirow{3}{*}{German Credit ($p\!=\!21$)} & CLT CI & 46.7 & \textbf{0.088} \\
 & BMCS & \textbf{100.0} & 0.188 \\
 & CBCS & --- & --- \\
\midrule
\multirow{3}{*}{Ionosphere ($p\!=\!34$)} & CLT CI & 40.0 & \textbf{0.129} \\
 & BMCS & \textbf{80.0} & 0.404 \\
 & CBCS & --- & --- \\
\midrule
\multirow{3}{*}{Pima ($p\!=\!9$)} & CLT CI & 86.7 & \textbf{0.058} \\
 & BMCS & \textbf{100.0} & 0.114 \\
 & CBCS & 0.0 & 0.765 \\
\bottomrule
\end{tabular}
\end{table}

\begin{figure}[t]
\centering
\includegraphics[width=\textwidth]{figures/fig3_logistic_regression.pdf}
\caption{Mean confidence interval/sequence width at the final iteration for each dataset. CLT intervals (gray) are narrowest but have poor sequential coverage. BMCS (blue) and MDCS (green) provide wider but more reliable intervals. Numbers above bars show time-uniform coverage.}
\label{fig:logistic}
\end{figure}

\subsection{Experiment 3: Robustness Under Poor Mixing}

\paragraph{Setup.} We target a mixture $0.5\N(-\mu_{\text{sep}}\mathbf{1}, I_5) + 0.5\N(\mu_{\text{sep}}\mathbf{1}, I_5)$ with true mean zero, varying the separation $\mu_{\text{sep}} \in \{1.0, 2.0, 3.0, 4.0\}$. We run 100 RWMH replicates per seed (300 total) with chains of length $100{,}000$.

\paragraph{Results.} Figure~\ref{fig:multimodal} and Table~\ref{tab:multimodal} show the results. At mild separation ($\mu_{\text{sep}} = 1.0$), the chain mixes reasonably (mean 2425 mode switches) and BMCS achieves 91.7\% time-uniform coverage---below nominal but far better than CLT's 61.3\%. At moderate separation ($\mu_{\text{sep}} = 2.0$), mixing deteriorates sharply (460 switches, IAT $\approx 1208$), and \emph{all} methods fail: BMCS drops to 2.0\% and CLT to 0\%. At $\mu_{\text{sep}} \geq 3.0$, the chain essentially traps in one mode ($<35$ switches) and no method provides valid coverage.

\begin{table}[t]
\caption{Time-uniform coverage (\%) on multimodal targets ($d=5$). Mode switches and integrated autocorrelation time (IAT) characterize mixing quality. All methods fail when mixing is poor ($\mu_{\text{sep}} \geq 2.0$).}
\label{tab:multimodal}
\centering
\begin{tabular}{lcccccc}
\toprule
$\mu_{\text{sep}}$ & Switches & IAT & CLT & BMCS & MDCS \\
\midrule
1.0 & 2425 & 67.8 & 61.3 & \textbf{91.7} & 85.3 \\
2.0 & 460 & 1207.9 & 0.0 & 2.0 & 3.0 \\
3.0 & 34 & 23.2 & 0.0 & 0.0 & 0.0 \\
4.0 & 1 & 18.5 & 0.0 & 0.0 & 0.0 \\
\bottomrule
\end{tabular}
\end{table}

\begin{figure}[t]
\centering
\includegraphics[width=\textwidth]{figures/fig4_multimodal_robustness.pdf}
\caption{Left: Time-uniform coverage vs.\ mode separation. All methods degrade as mixing worsens, but BMCS (blue) degrades more gracefully than CLT (gray) at mild separation. Right: Mixing diagnostics showing mode switches (bars) and integrated autocorrelation time (red line).}
\label{fig:multimodal}
\end{figure}

This experiment reveals an important limitation: \emph{confidence sequences cannot substitute for chain convergence}. When the MCMC chain fails to explore the full target distribution, no post-hoc statistical method can recover valid coverage. The confidence sequence width reflects the variability observed within the trapped mode, not the true posterior uncertainty.

\subsection{Experiment 4: Stopping Efficiency}

\paragraph{Setup.} We compare stopping strategies on Gaussian targets ($d \in \{1, 10, 50\}$) with target half-widths $\epsilon \in \{0.1, 0.05\}$. For each of 80 replicates per seed (240 total), we compute the stopping time for FWSR (CLT-based) and BMCS, and compare against the oracle stopping time $n^* = z_{\alpha/2}^2 \sigma^2_\infty / \epsilon^2$.

\paragraph{Results.} Table~\ref{tab:stopping} presents the stopping efficiency comparison. A key finding is that \textbf{BMCS requires more samples than FWSR}, not fewer. BMCS stopping times are $1.7$--$3.4\times$ the oracle, while FWSR stopping times are $0.58$--$0.99\times$ the oracle. However, this comparison reveals that \textbf{FWSR is anti-conservative}: FWSR achieves coverage of only 87.9\%--94.6\% (below the 95\% nominal) in most settings, sometimes reaching 0\% when the chain length is insufficient. BMCS, by contrast, achieves $\geq 98.8\%$ coverage whenever it stops (though it sometimes does not stop within the maximum chain length).

\begin{table}[t]
\caption{Stopping efficiency. Ratio = stopping time / oracle stopping time ($\downarrow$). Coverage is at the stopping time ($\uparrow$). FWSR stops earlier but with invalid coverage; BMCS stops later but with valid coverage. ``---'' means method did not stop within chain length.}
\label{tab:stopping}
\centering
\begin{tabular}{llcccc}
\toprule
Target & $\epsilon$ & \multicolumn{2}{c}{FWSR} & \multicolumn{2}{c}{BMCS} \\
\cmidrule(lr){3-4} \cmidrule(lr){5-6}
 & & Ratio $\downarrow$ & Cov.\ (\%) $\uparrow$ & Ratio $\downarrow$ & Cov.\ (\%) $\uparrow$ \\
\midrule
Gauss $d\!=\!1$ & 0.10 & 0.89 & 91.7 & 2.94 & 98.8 \\
Gauss $d\!=\!1$ & 0.05 & 0.89 & 94.6 & 3.36 & 99.6 \\
Gauss $d\!=\!10$ & 0.10 & 0.77 & 87.9 & 3.16 & 100.0 \\
Gauss $d\!=\!10$ & 0.05 & 0.86 & 92.9 & --- & --- \\
Gauss $d\!=\!50$ & 0.10 & 0.58 & 89.6 & --- & --- \\
Gauss $d\!=\!50$ & 0.05 & 1.00 & 0.0 & --- & --- \\
\bottomrule
\end{tabular}
\end{table}

\begin{figure}[t]
\centering
\includegraphics[width=0.6\textwidth]{figures/fig5_stopping_efficiency.pdf}
\caption{Stopping time relative to oracle for FWSR and BMCS across settings. FWSR (red) consistently stops near or below the oracle but with invalid coverage. BMCS (blue) requires $2$--$3\times$ more samples but maintains valid coverage.}
\label{fig:stopping}
\end{figure}

The stopping results reveal that the ``price of anytime validity'' manifests as requiring more samples, not less. This is the expected cost of maintaining coverage at all intermediate stopping times. The original hypothesis that confidence sequences would enable computational \emph{savings} is not supported; instead, the value is in providing \emph{valid} guarantees when practitioners do choose to stop early.

\subsection{Experiment 5: Scalability}

\paragraph{Setup.} We evaluate computational overhead and width scaling for $d \in \{2, 5, 10, 20, 50, 100, 200\}$ on Gaussian targets with $n = 50{,}000$ and 50 replicates per seed.

\paragraph{Results.} Figure~\ref{fig:scalability} shows that BMCS adds very low computational overhead: the ratio of BMCS computation time to chain running time decreases from 20\% at $d=2$ to 5.6\% at $d=200$, because the MCMC chain's per-iteration cost grows with~$d$ while the BMCS bookkeeping remains constant. The BMCS width is approximately $1.9$--$2.9\times$ the CLT width across dimensions. CBCS becomes impractical for $d \geq 20$ as the meeting time reaches the cap of 5000 iterations.

\begin{table}[t]
\caption{Scalability results. Chain time is per-iteration cost ($\mu$s). Overhead ratio is CS computation time / chain time. Width ratio is BMCS width / CLT width. Meeting time is the median for CBCS (capped at 5000).}
\label{tab:scalability}
\centering
\begin{tabular}{rcccccc}
\toprule
$d$ & Chain ($\mu$s) & BMCS Ovhd.\ & Width (CLT) $\downarrow$ & Width (BMCS) $\downarrow$ & Ratio & CBCS $\tau$ \\
\midrule
2 & 12.0 & 20.4\% & 0.025 & 0.045 & 1.84 & 5 \\
5 & 12.7 & 20.1\% & 0.035 & 0.065 & 1.85 & 15 \\
10 & 13.1 & 19.5\% & 0.048 & 0.089 & 1.86 & 213 \\
20 & 21.6 & 11.9\% & 0.064 & 0.124 & 1.93 & 5000 \\
50 & 23.7 & 10.8\% & 0.089 & 0.189 & 2.13 & 5000 \\
100 & 28.5 & 9.0\% & 0.105 & 0.257 & 2.45 & 5000 \\
200 & 46.2 & 5.6\% & 0.115 & 0.334 & 2.90 & 5000 \\
\bottomrule
\end{tabular}
\end{table}

\begin{figure}[t]
\centering
\includegraphics[width=\textwidth]{figures/fig6_scalability.pdf}
\caption{Left: Computational overhead ratio (CS time / chain time) vs.\ dimension. BMCS (blue) adds $<20\%$ overhead at all dimensions. Right: Confidence interval/sequence width at $n=50{,}000$ vs.\ dimension. BMCS (blue) is $\sim 2\times$ wider than CLT (gray). MDCS (green) produces the widest intervals.}
\label{fig:scalability}
\end{figure}

\subsection{Ablation Studies}

We conduct ablation studies to evaluate the contribution of each BMCS component. Table~\ref{tab:ablation} summarizes results on a Gaussian target with $d=10$ and 200 replicates per seed (600 total).

\begin{table}[t]
\caption{Ablation studies for BMCS on Gaussian $d=10$. Time-uniform coverage (\%) and mean width at $n=50{,}000$. 600 replicates per configuration.}
\label{tab:ablation}
\centering
\begin{tabular}{llcc}
\toprule
Ablation & Variant & Coverage (\%) $\uparrow$ & Width $\downarrow$ \\
\midrule
\multirow{3}{*}{A1: Batch size} & Adaptive (ours) & \textbf{98.7} & 0.089 \\
 & $\sqrt{n}$ & 94.3 & 0.109 \\
 & $n^{1/3}$ & 98.0 & \textbf{0.089} \\
\midrule
\multirow{2}{*}{A2: Dependence corr.} & With correction & \textbf{98.7} & 0.089 \\
 & Without correction & 98.8 & \textbf{0.095} \\
\midrule
\multirow{2}{*}{A3: CS type} & Emp.\ Bernstein (ours) & \textbf{98.7} & \textbf{0.089} \\
 & Hoeffding & 100.0 & 0.157 \\
\midrule
\multirow{3}{*}{A4: Burn-in} & No burn-in (ours) & \textbf{98.7} & \textbf{0.089} \\
 & 10\% discard & 97.5 & 0.141 \\
 & 20\% discard & 98.3 & 0.141 \\
\bottomrule
\end{tabular}
\end{table}

\begin{figure}[t]
\centering
\includegraphics[width=\textwidth]{figures/fig7_ablation.pdf}
\caption{Time-uniform coverage across ablation variants. All BMCS variants maintain coverage near or above 95\% on Gaussian targets. The multimodal target (A2, center) shows that all methods fail under poor mixing regardless of the correction term.}
\label{fig:ablation}
\end{figure}

\textbf{A1 (Batch size):} Our adaptive strategy achieves the best coverage (98.7\%) with competitive width. The $\sqrt{n}$ strategy produces too few batches early on, degrading coverage to 94.3\%.

\textbf{A2 (Dependence correction):} On well-mixing Gaussian targets, the correction has minimal effect (98.7\% vs.\ 98.8\%). On the multimodal target ($\mu_{\text{sep}} = 2.0$), neither variant achieves valid coverage ($\sim$5\%), confirming that the correction cannot compensate for fundamental mixing failure.

\textbf{A3 (CS type):} The empirical Bernstein CS provides both good coverage (98.7\%) and tight width (0.089). The Hoeffding CS achieves 100\% coverage but is $1.76\times$ wider (0.157), as it does not adapt to the empirical variance.

\textbf{A4 (Burn-in):} Discarding burn-in samples is unnecessary and harmful: it reduces the effective number of batches without improving coverage, increasing the final width by $\sim 58\%$ (0.089 to 0.141).

\paragraph{CBCS ablations.} For CBCS on Gaussian $d=10$ with 100 coupled pairs per seed: (i) maximal coupling achieves median meeting time $\tau = 193$ while common random numbers (CRN) coupling fails entirely ($\tau = 5000$ cap); (ii) the burn-in parameter $k = \lfloor 3\tau/4 \rfloor$ yields the lowest variance estimators (std = 8.18 vs.\ 20.27 for $k = \lfloor \tau/2 \rfloor$); (iii) for equal computational budget, BMCS achieves width 0.050 vs.\ CBCS width 7.57---BMCS is over $150\times$ more sample-efficient.

% ============================================================================
% DISCUSSION
% ============================================================================
\section{Discussion and Limitations}
\label{sec:discussion}

\paragraph{Key findings.} Our experiments establish that confidence sequences for MCMC are practical and provide substantially stronger guarantees than CLT intervals under sequential monitoring. BMCS is the recommended default: it maintains $\geq 93.7\%$ time-uniform coverage on well-mixing chains, adds less than 20\% computational overhead, and requires no coupling infrastructure. The price of anytime validity is approximately $1.9\times$ wider intervals and, when used as a stopping rule, $2$--$3\times$ more samples than FWSR---but FWSR's coverage is invalid under the sequential monitoring it implicitly assumes.

\paragraph{CBCS limitations.} CBCS provides exact independence but is impractical in most settings. Meeting times grow rapidly with dimension, becoming infeasible for $d \geq 20$. Even when coupling succeeds, the per-estimator variance is high, leading to wide intervals that are dominated by BMCS at equal computational cost. We include CBCS for theoretical completeness but do not recommend it for routine use.

\paragraph{MDCS limitations.} The martingale decomposition approach is theoretically elegant but practically unreliable. Accurate estimation of the Poisson equation solution norm is difficult, especially in high dimensions, leading to unreliable coverage (as low as 65.7\% at $d=50$). Future work on better Poisson equation solvers could make this method competitive.

\paragraph{Mixing failure.} No confidence sequence method can compensate for poor MCMC mixing. When the chain fails to explore all modes (as in our multimodal experiments with $\mu_{\text{sep}} \geq 2.0$), all methods fail catastrophically. Confidence sequences should be viewed as tools for valid inference \emph{conditional on adequate mixing}, not as mixing diagnostics.

\paragraph{Width overhead.} The BMCS width ratio to CLT increases with dimension ($1.84\times$ at $d=2$ to $2.90\times$ at $d=200$). For applications requiring very precise estimates, this overhead may be substantial. Understanding whether tighter confidence sequences are achievable for MCMC is an important open question.

\paragraph{Stopping efficiency.} Contrary to our initial hypothesis, BMCS does not enable earlier stopping than FWSR; it requires more samples. The value proposition is not computational savings but \emph{inferential validity}: BMCS provides guaranteed coverage at any stopping time, while FWSR's coverage guarantees are invalid when used as intended (sequential monitoring until the interval is narrow enough).

% ============================================================================
% CONCLUSION
% ============================================================================
\section{Conclusion}

We developed the first confidence sequences for MCMC estimators, bridging the gap between confidence sequence theory and Bayesian computation. Our Batch Means Confidence Sequence (BMCS) method maintains $\geq 93.7\%$ time-uniform coverage on Gaussian targets up to $d=50$ while adding less than 20\% computational overhead. In contrast, standard CLT intervals degrade to as low as 28.3\% coverage under sequential monitoring---the very setting in which they are universally used.

Our experiments reveal both the value and the cost of anytime validity for MCMC. The value is clear: valid coverage guarantees at every stopping time, eliminating the coverage violations inherent in current practice. The cost is also clear: approximately $2\times$ wider intervals and, when used as stopping rules, $2$--$3\times$ more computation. We argue that this tradeoff favors validity, as the CLT interval's nominal coverage is simply not what practitioners actually receive.

Important limitations remain. BMCS coverage degrades in high dimensions ($d=50$) and when chains mix poorly. The coupling-based CBCS method, while theoretically exact, is impractical beyond low dimensions. Future work should develop tighter confidence sequences for MCMC, investigate the theoretical achievable rate in the dependent setting, and explore adaptive methods that simultaneously diagnose mixing failure and provide valid estimation.

% ============================================================================
% REFERENCES
% ============================================================================
\bibliographystyle{plainnat}

\begin{thebibliography}{15}

\bibitem[Corenflos and Dau(2025)]{corenflos2025coupling}
Adrien Corenflos and Hai-Dang Dau.
\newblock A coupling-based approach to f-divergences diagnostics for {M}arkov chain {M}onte {C}arlo.
\newblock \emph{arXiv preprint arXiv:2510.07559}, 2025.

\bibitem[Flegal and Jones(2010)]{flegal2010batch}
James~M. Flegal and Galin~L. Jones.
\newblock Batch means and spectral variance estimators in {M}arkov chain {M}onte {C}arlo.
\newblock \emph{The Annals of Statistics}, 38(2):1034--1070, 2010.

\bibitem[Flegal et~al.(2024)Flegal, Hughes, Vats, and Dai]{flegal2024implementing}
James~M. Flegal, James Hughes, Dootika Vats, and Ning Dai.
\newblock Implementing {MCMC}: Multivariate estimation with confidence.
\newblock \emph{arXiv preprint arXiv:2408.15396}, 2024.

\bibitem[Geyer(1992)]{geyer1992practical}
Charles~J. Geyer.
\newblock Practical {M}arkov chain {M}onte {C}arlo.
\newblock \emph{Statistical Science}, 7(4):473--483, 1992.

\bibitem[Geyer(2011)]{geyer2011introduction}
Charles~J. Geyer.
\newblock Introduction to {M}arkov chain {M}onte {C}arlo.
\newblock In Steve Brooks, Andrew Gelman, Galin~L. Jones, and Xiao-Li Meng, editors, \emph{Handbook of Markov Chain Monte Carlo}, chapter~1. Chapman and Hall/CRC, 2011.

\bibitem[Howard et~al.(2021)Howard, Ramdas, McAuliffe, and Sekhon]{howard2021time}
Steven~R. Howard, Aaditya Ramdas, Jon McAuliffe, and Jasjeet Sekhon.
\newblock Time-uniform, nonparametric, nonasymptotic confidence sequences.
\newblock \emph{The Annals of Statistics}, 49(2):1--36, 2021.

\bibitem[Jacob et~al.(2020)Jacob, O'Leary, and Atchad\'{e}]{jacob2020unbiased}
Pierre~E. Jacob, John O'Leary, and Yves~F. Atchad\'{e}.
\newblock Unbiased {M}arkov chain {M}onte {C}arlo methods with couplings.
\newblock \emph{Journal of the Royal Statistical Society: Series B}, 82(3):543--600, 2020.

\bibitem[Martinez-Taboada and Ramdas(2025)]{martinez2025sequential}
Diego Martinez-Taboada and Aaditya Ramdas.
\newblock Sequential kernelized {S}tein discrepancy.
\newblock In \emph{Proceedings of AISTATS}, 2025.

\bibitem[Ramdas et~al.(2023)Ramdas, Gr\"{u}nwald, Vovk, and Shafer]{ramdas2023game}
Aaditya Ramdas, Peter Gr\"{u}nwald, Vladimir Vovk, and Glenn Shafer.
\newblock Game-theoretic statistics and safe anytime-valid inference.
\newblock \emph{Statistical Science}, 38(4):576--601, 2023.

\bibitem[Robertson et~al.(2021)Robertson, Flegal, Vats, and Jones]{robertson2021assessing}
Nathan Robertson, James~M. Flegal, Dootika Vats, and Galin~L. Jones.
\newblock Assessing and visualizing simultaneous simulation error.
\newblock \emph{Journal of Computational and Graphical Statistics}, 30(4):1157--1168, 2021.

\bibitem[Vats and Flegal(2022)]{vats2022lugsail}
Dootika Vats and James~M. Flegal.
\newblock Lugsail lag windows for estimating time-average covariance matrices.
\newblock \emph{Biometrika}, 109(3):735--750, 2022.

\bibitem[Vats et~al.(2019)Vats, Flegal, and Jones]{vats2019multivariate}
Dootika Vats, James~M. Flegal, and Galin~L. Jones.
\newblock Multivariate output analysis for {M}arkov chain {M}onte {C}arlo.
\newblock \emph{Biometrika}, 106(2):321--337, 2019.

\bibitem[Wang and Ramdas(2023)]{wang2023catoni}
Hongjian Wang and Aaditya Ramdas.
\newblock Catoni-style confidence sequences for heavy-tailed mean estimation.
\newblock \emph{Stochastic Processes and their Applications}, 163:51--79, 2023.

\end{thebibliography}

\end{document}
